\documentclass[aspectratio=169,11pt]{beamer}
\def\courseweek{8}
% Shared style for the Introduction to Information Geometry lecture slides.
% Compile the lecture .tex files from this directory.
\providecommand{\courseweek}{1}
\usepackage[T1]{fontenc}
\usepackage[utf8]{inputenc}
\usepackage[sc]{mathpazo}
\renewcommand{\sfdefault}{lmss}
\renewcommand{\ttdefault}{lmtt}
\usepackage{amsmath,amssymb,mathtools}
\usepackage{tikz}
\usetikzlibrary{arrows.meta,calc,positioning}
\usepackage{pgfplots}
\pgfplotsset{compat=1.18}
\usepackage{booktabs}

\definecolor{courseink}{HTML}{173749}
\definecolor{coursemuted}{HTML}{526674}
\definecolor{courseteal}{HTML}{0E7287}
\definecolor{coursecopper}{HTML}{B66B43}
\definecolor{courseline}{HTML}{D4E0E3}
\definecolor{coursepale}{HTML}{EFF5F5}

\setbeamercolor{normal text}{fg=courseink,bg=white}
\setbeamercolor{structure}{fg=courseteal}
\setbeamercolor{frametitle}{fg=courseink,bg=white}
\setbeamercolor{title}{fg=courseink}
\setbeamercolor{subtitle}{fg=coursemuted}
\setbeamercolor{block title}{fg=courseink,bg=white}
\setbeamercolor{block body}{fg=courseink,bg=white}
\setbeamercolor{alerted text}{fg=coursecopper}
\setbeamerfont{title}{family=\rmfamily,series=\bfseries,size=\fontsize{32}{37}}
\setbeamerfont{frametitle}{family=\rmfamily,series=\bfseries,size=\fontsize{26}{31}}
\setbeamerfont{normal text}{size=\fontsize{18}{22}}
\setbeamerfont{subtitle}{size=\fontsize{17}{21}}
\setbeamersize{text margin left=0.78cm,text margin right=0.78cm}
\setbeamertemplate{navigation symbols}{}
\setbeamertemplate{itemize items}[circle]
\setbeamertemplate{enumerate items}[default]
\setbeamertemplate{blocks}[default]
\setlength{\parskip}{0.22em}
\setbeamertemplate{footline}{%
  \leavevmode\hbox to \paperwidth{\hspace{.78cm}%
  {\color{coursemuted}\fontsize{8.5}{10}\selectfont Introduction to Information Geometry \hspace{.32cm} / \hspace{.32cm} Week \courseweek}%
  \hfill{\color{coursemuted}\fontsize{8.5}{10}\selectfont \insertframenumber\hspace{.17cm}/\hspace{.17cm}\inserttotalframenumber}\hspace{.78cm}}%
  \vskip.17cm
}
\setbeamertemplate{frametitle}{%
  \vspace{.13cm}
  \insertframetitle\par
  \vspace{.08cm}
}
\setbeamertemplate{title page}{%
  \vspace*{1.00cm}
  {\color{courseteal}\fontsize{10}{13}\selectfont\bfseries
  INTRODUCTION TO INFORMATION GEOMETRY\par}
  \vspace{.40cm}
  {\usebeamerfont{title}\usebeamercolor[fg]{title}\inserttitle\par}
  \vspace{.22cm}
  {\usebeamerfont{subtitle}\usebeamercolor[fg]{subtitle}\insertsubtitle\par}
  \vfill
  {\color{coursemuted}\fontsize{13}{16}\selectfont Christopher Lee\par}
  \vspace{.34cm}
}

\pgfplotsset{
  courseplot/.style={
    width=11.05cm,
    height=5.15cm,
    axis line style={coursemuted!75},
    tick style={coursemuted!65},
    ticklabel style={font=\scriptsize,color=coursemuted},
    label style={font=\normalsize,color=courseink},
    grid=major,
    major grid style={courseline!70,line width=.4pt},
    line width=1.45pt,
    legend style={font=\small,draw=none,fill=none,at={(.98,.03)},anchor=south east},
    clip=true
  }
}

\newcommand{\highlightmath}[1]{{\color{courseteal}#1}}
\newcommand{\thinrule}{\par\vspace{.25cm}{\color{courseline}\rule{\linewidth}{.8pt}}\par}

% A one-dimensional state space.  x=0 and x=1 are placed at 0 and 9 units.
\newcommand{\stateinterval}[1]{%
 \begin{tikzpicture}[x=1cm,y=1cm]
   \draw[courseink,line width=1.25pt] (0,0)--(9,0);
   \foreach \pos/\lab in {0/0,9/1} {
     \fill[courseink] (\pos,0) circle(2.7pt);
     \node[below=9pt,text=coursemuted] at (\pos,0) {\(\lab\)};
   }
   #1
 \end{tikzpicture}%
}

% Paired pages are separate frames in the PDF, with matching titles.
% Supply the whole frame body for each version; the instructor may edit them
% independently without changing overlay logic.
\newcommand{\coursepair}[3]{%
  \begin{frame}{#1}#2\end{frame}%
  \begin{frame}{#1}#3\end{frame}%
}

% Equilateral representation of the categorical simplex, in barycentric
% coordinates with vertices A=(0,0), B=(1,0), C=(1/2,sqrt(3)/2).
\newcommand{\coursesimplex}[1]{%
 \begin{tikzpicture}[scale=4.9]
  \fill[coursepale] (0,0)--(1,0)--(.5,.8660254)--cycle;
  \draw[courseink,line width=.65pt] (0,0)--(1,0)--(.5,.8660254)--cycle;
  \node[below left,font=\small] at (0,0) {A};
  \node[below right,font=\small] at (1,0) {B};
  \node[above,font=\small] at (.5,.8660254) {C};
  #1
 \end{tikzpicture}%
}

\title{Geodesics and Statistical Midpoints}
\subtitle{Week 8 \(\cdot\) Lecture 2}
\date{}
\begin{document}
\begin{frame}[plain]\titlepage\end{frame}

\begin{frame}{A constant-speed spherical geodesic}
Map \(P\) and \(R\) to \(s_0=2\sqrt P\) and \(s_1=2\sqrt R\).
Let \(\alpha=\arccos B(P,R)\).
\vspace{.40cm}
Then
\[
\highlightmath{s(t)=\frac{\sin((1-t)\alpha)}{\sin\alpha}s_0
+\frac{\sin(t\alpha)}{\sin\alpha}s_1}.
\]
\vspace{.38cm}
The probability path is \(\gamma_i(t)=s_i(t)^2/4\).
\end{frame}

\begin{frame}{Why the speed is constant}
Choose an orthonormal basis \(e_0,e_1\) for the plane of the great circle.
Then the same path has the form
\[
s(t)=2\big(\cos(t\alpha)e_0+\sin(t\alpha)e_1\big).
\]
\vspace{.50cm}
Differentiate:
\[
\highlightmath{\|s'(t)\|=2\alpha}.
\]
\vspace{.42cm}
Its length over \(0\leq t\leq1\) is \(2\alpha=d_F(P,R)\).
\end{frame}

\begin{frame}{The Fisher midpoint formula}
At \(t=\frac12\), both sine coefficients agree.
The probability midpoint between \(P\) and \(R\) is
\vspace{.4cm}
\[
\highlightmath{\gamma_i(\tfrac12)=
\frac{(\sqrt{p_i}+\sqrt{r_i})^2}{2(1+B(P,R))}}.
\]
\vspace{.42cm}
The squares make every coordinate nonnegative;
the denominator makes them sum to one.
\end{frame}

\coursepair{The Bernoulli midpoint}{
Find the constant-speed Fisher midpoint between \(p=\frac14\)
and \(p=\frac34\).
\vspace{.43cm}
Recall \(\theta(p)=2\arcsin\sqrt p\).
\vspace{.42cm}
The midpoint is \(p_{\mathrm{mid}}=\;?\).
}{
Find the constant-speed Fisher midpoint between \(p=\frac14\)
and \(p=\frac34\).
\vspace{.43cm}
The arc coordinates are \(\pi/3\) and \(2\pi/3\).
\vspace{.42cm}
Their midpoint is \(\pi/2\), so \(\highlightmath{p_{\mathrm{mid}}=\frac12}\).
}

\begin{frame}{Mixtures trace straight probability segments}
The ordinary mixture between distributions is
\[
M(t)=(1-t)P+tR.
\]
\vspace{.48cm}
It is a straight segment in the affine simplex drawing.
\vspace{.37cm}
The Fisher geodesic comes from a great-circle arc after taking
square roots. For typical three-outcome endpoints, the two paths differ.
\end{frame}

\begin{frame}{Two paths between the same laws}
\begin{center}\scalebox{.92}{\coursesimplex{
 \coordinate (p) at (.22,.13);
 \coordinate (r) at (.69,.45);
 \draw[coursecopper,line width=1.3pt](p)--(r);
 \draw[courseteal,line width=1.8pt](p)..controls(.30,.47)and(.55,.57)..(r);
 \fill[courseink](p)circle(2pt);\fill[courseink](r)circle(2pt);
 \node[below left,font=\small]at(p){\(P\)};
 \node[above right,font=\small]at(r){\(R\)};
}}\end{center}
\vspace{.1cm}
Schematic: copper shows a mixture segment; teal shows a curved Fisher route.
\end{frame}

\begin{frame}{Hellinger distance and spherical chords}
The square-root overlap also determines a Euclidean chord:
\[
\|2\sqrt P-2\sqrt R\|^2=8(1-B(P,R)).
\]
\vspace{.42cm}
In one common normalization, Hellinger distance satisfies
\[
\highlightmath{h(P,R)^2=1-B(P,R)}.
\]
\vspace{.38cm}
The chord and the spherical arc agree to first order for nearby laws.
\end{frame}

\coursepair{Chord and arc between pure laws}{
Take \(P=(1,0)\) and \(R=(0,1)\).
\vspace{.47cm}
Their overlap is \(B=0\). Find
\[
h(P,R),\quad \|2\sqrt P-2\sqrt R\|,\quad d_F(P,R).
\]
}{
Take \(P=(1,0)\) and \(R=(0,1)\).
\vspace{.47cm}
Their overlap is \(B=0\), so
\[
\highlightmath{h(P,R)=1},\qquad
\highlightmath{\|2\sqrt P-2\sqrt R\|=2\sqrt2},
\]
\vspace{.18cm}
\[
\highlightmath{d_F(P,R)=\pi}.
\]
}

\begin{frame}{Global distance and the local metric agree}
For a small tangent change \(h\), with \(\sum_i h_i=0\),
\[
B(P,P+h)=1-\frac18\sum_i\frac{h_i^2}{p_i}
+O(\|h\|^3).
\]
\vspace{.43cm}
It follows that
\[
\highlightmath{d_F(P,P+h)^2=\sum_i\frac{h_i^2}{p_i}
+O(\|h\|^3)}.
\]
\end{frame}

\begin{frame}{Curvature and the simplex boundary}
The categorical Fisher manifold is an open region of
the radius-2 sphere.
\vspace{.45cm}
Its sectional curvature is
\[
\highlightmath{K=\frac1{2^2}=\frac14}.
\]
\vspace{.43cm}
The positive interior is incomplete. Including distributions with zero
coordinates gives its metric completion as the closed spherical orthant.
\end{frame}

\begin{frame}{Two geometries of the same state space}
The affine simplex represents mixtures as straight segments.
\vspace{.48cm}
The Fisher metric represents the same probability states
as an orthant of a sphere.
\vspace{.48cm}
Next week we will use that metric to construct gradient flows
on spaces of distributions.
\end{frame}
\end{document}
